\chapter{琼楼玉宇}

“原来，向下坠落的是玩偶，而向上坠落的是我。”玲望着面前这片包裹着她的猩红的夜空；那些原本遥不可及的星辰，此刻却像是被无形之手拉近，闪烁着诡异的十字形光芒，飞快地向她逼近。这景象瓦解了她对宇宙尺度的所有认知。“这下我也‘手可摘星辰’了。”玲合时宜地想起李白的诗，“这并不奇怪；毕竟，‘危楼’只需要‘高百尺’，也就是差不多三十米；而刚才的楼显然超过十层了。”如果说李白的诗表达了建筑的高度和摘星星的能力之间的因果关系，那么这样一种因果关系在此时此刻成立了。“不过，”玲又想，“如果让我来写，我肯定写‘危楼高千尺’；毕竟‘百尺’怎么看都不算很高。”然而玲立刻否定了自己的观点，“‘百’是仄声，‘千’是平声；如果这样改，就不符合音韵了。”玲于是开始思考若是把这句放到第三句，她能够进行怎样的拗救。意识到自己想不出主意后，玲叹了口气，感到自己的理性正在这片血色虚空中无谓地空转，“要是李白写的是阿拉伯语的诗就好了；这样，‘千’和‘夜’就刚好能组成一对，用来描写夜空下的千尺高楼。”李白出生于碎叶城；阿拉伯商旅往来；李白会点儿阿拉伯语也不奇怪吧？正思考着，玲突然灵光一现，“直接改成‘万尺’不就行了；‘万’也是仄声。这下应该足够高了；现在以及未来一段时间的工程学应该无法建造这么高的建筑。”玲于是开始想象万尺高的碎叶城（这没有造成事实谬误；因为李白没有描写过他的出生之地；既然如此，任何想象便是自由的）。脱落了几个字之后，玲便想到了《万叶集》。“岛夷真可怜，到了差不多唐朝的时候才写出《诗经》的类似物。”不过，玲这样的想法包含侮辱性称呼，显然是不能公开的。“怎么不能公开；又不是专指；《魏书》不也骂南朝为岛夷。”玲回忆着《万叶集》中的短歌，发现自己一首都不记得了。“开头都是五、七、五；我现在不如来写俳句吧！”

于是玲开始写俳句：
\begin{quote}
猩红色天空\\
向上坠落每一刻\\
都觉得无聊
\end{quote}

这样是不行的；总不能拿“猩红色天空”当季语吧！那将会是一个怎样的季节。于是玲搜肠刮肚，正经地写了一首：
\begin{quote}
冻星垂血穹\\
苇笛裂帛断夜半\\
独钓寒江雪
\end{quote}

好一个“独钓寒江雪”！玲为自己跳跃的思路感到遗憾，并对自己目前的境况作出结论，“要是继续这样无止境地坠落，给我一台打字机，我都能写出莎士比亚的全部著作了。”基于这样的观点，我们可能并不需要先进的AI模型，而仅仅需要一只猴子和无限的时间。不过，对于坠落时间的简单外推被证明是不合理的；因为一会儿之后，那无声的、永恒般的坠落突兀地终止了。某颗原本只是光点的星星在她眼前急剧放大，从血色天穹上垂下来，露出反射着诡异光泽的表面。玲轻轻地落到其上，脚下的星辰传递来一种奇特的反馈感。透过鞋底，玲能感觉到地面异常致密和平滑，同时又带着一种难以言喻的、隐约的温和感，仿佛这巨大的平面本身就是一个低功率的热源。这种感觉如此独特，让她联想到行走在微温的、巨大骨瓷工艺品上的荒谬景象。

远处是一个建筑群的朦胧的轮廓。玲跑过去（今天的运动量已经超标了），并在跑动的过程中判定该星球的重力和地球是相当的。说不定这就是地球，而远处就是一个普通的市镇；那样的话，玲就可以搭便车回家睡觉了。然而，靠近的观察让玲失望了：眼前的景象彻底背弃了所有关于建筑和结构力学的常识。巨大的、仿佛被神祇随意掷下的倾斜四棱柱构成了这里的一切。它们以匪夷所思的角度互相支撑、穿插，平滑的切面则反射着一种荧白色的光芒。说不定这就是现代建筑艺术。玲上前，想要知道这些四棱柱的材质。她摸了摸；石块有种温润而细腻的触感，就像玉一样。

“嘿，这就是玉！”玲终于反应过来，“玉的矿物组分是单斜晶系的，确实就该是倾斜的四棱柱。”她环顾四周，这哪里是建筑群，分明是一座完全由巨型玉石单晶构成的山脉，一座被强行赋予了宫殿形态的地质奇观。这么多这么大块的玉的单晶体，肯定很值钱。要是手上有把地质锤，能够凿点下来带回去就好了。（事实上小块的角闪石并不值钱，而且还能人工合成；当然，如果玲能够搬回去一整块，还是很有经济利益的。）

“如果这里都是些玉制建筑，”玲重新开始思考自己的处境，“那么它们一定就是苏轼所见的‘琼楼玉宇’了。”于是玲确定自己就像爱丽丝一样来到了仙境，“这是中国读者因译名而产生的一个常见的误解：刘易斯所写的Wonderland，说的是奇境，而并不专指仙境。”玲赶紧纠正自己。

玲顺着这些围绕起来的巨大晶块走着，脚下的每一步都回荡着清脆而空旷的声响；她想要找到通往这个建筑群内部的入口。她来到这个建筑群（现在应该叫它“天宫”了）的正南方。面前是四根巨大的直插天际的柱子（当然也是倾斜的），柱子下看起来是一些执勤人员。那么这就是门了；但是这门怎么没有门楣呢？或许是在无穷高的天上吧；天宫的天上又是什么呢？（如果玲亲身观察过古建筑的话，她会认识到这是华表，并不是门。）

“站住！”看到玲接近，执勤人员探过头来恶狠狠地呵道。

“探过头来”的意思是他身体不动而把头从十几米开外伸过来了。玲发现这位高大（横向）的执勤人员的脖子竟在不停地伸长，弯弯绕绕而随意地在地上盘了好几圈。

“你好，我要登记参观。”

“那你得先填表。”

“那你得先给我表。”

“表不在这儿领，”长脖子保安不耐烦地说着，脖颈蠕动了一下，带动地上的圈儿也随之起伏，“你去找访客管理处。”

“访客管理处在哪里？”

“在里面。”

好吧，这个卡夫卡式的笑话一点儿也不好笑。可能的原因在于，玲心想，其一，这个模式过于老套了；其二，本书的作者并没有卡夫卡的语言功力。接下来的套路就应该是主角使用一些奇特的方法通过这项障碍了；不过玲目前还没有想到有什么方法。长脖子仍然盯着玲，瞪大的眼睛好似灯泡；玲只好尝试搭话以缓和这样的尴尬局面：

“你是谁？你的脖子为什么一直在变长？”玲很好奇这点；因为按照刘易斯的原著，应该是她的脖子在变长才对。但玲立刻后悔了；这个问题涉及对他人的生理缺陷的歧视。

长脖子显然不想搭理玲，而是想赶玲走；但是他还是回答了：“我是南天门守卫，增长天王。”他立刻补上一句：“不是‘脖子增长’的增长；是‘业绩增长’的增长。如果你念错了，我就斩断你的业绩。”他装腔作态地挥了挥手上的宝剑。

增长天王不管智慧来管业绩了，这可是个新鲜事。玲开始思考“业绩”是否是一个来自梵语的佛教词汇；最后她得出结论，尽管佛教中有“业”的概念，但“业绩”这词应该不来自佛教。“我会牢记的，”玲应付着，“所以您能通融一下吗？我只是想参观一下。”

“不行！”长脖子断然否决，“我有拒客率的指标。别妨碍我的业绩。”

要求拒客率“增长”也太奇怪了。这样的量标仅仅是拒绝的样本数量除以总的样本数量，而根本没有考虑“正确的拒绝”和“错误的拒绝”的区别。要是玲训练模型的时候能只被要求拒绝率增长就太好了。“医保公司肯定很喜欢这样的模型。”玲想到了一个用途。话说回来，这是服务业应有的态度吗？（保安确实属于服务业。）不对，服务业的态度是面向服务对象的；无许可的来者可不是保安的服务对象。

“任何人不得入内，”长脖子宣布着严肃的戒律，“除非他在列明的允许通行名单内。”他的头围绕着玲转了一圈，打量着玲，“你叫什么？”幸好玲是短发，不会被缠到。

“玲。”

“名单里没有这样的名字。”

“那，”玲忽然想赌一把，“你也可以叫我爱丽丝。”

“年轻人还起个英文名。”长脖子手忙脚乱地把一份长得同样拖到地上的名单往前翻。毕竟“爱丽丝”的首字母是“A”嘛！按道理是应该出现在名单的最开头一部分。玲心想，他这样翻名单，还不如直接把头伸到名单开头去看；不过，要是名单比长脖子的脖子还长，那么这个方法就不适用了。看来天庭是时候引入数字化管理了。长脖子检查了一会儿，念道，“爱丽丝，数据女工。”他又回过头打量了玲几眼，眼神中充满了对“数据女工”该有样貌的评判，“以后填表的时候不许写英文名了！你又不是国外的神仙机构派过来交流的。”

玲对于自己被当成“数据女工”非常不满；但现在显然不是抱怨的时候。“谢谢师傅！”玲完成基础的礼节，髂腰肌准备发力，想要赶紧进到门里去。“注意增长你的业绩，”不知道算不算好心的提醒，“我们这儿最不缺的就是数据女工。”

“我们确实应该提高数据标注员的待遇。”玲这样反思着，快步走进了那四根倾斜玉柱构成的、没有实体的天宫的大门。

天宫内部如同一个庞大的迷宫。玉石制成的重重楼阁倾斜而扭曲地交错在一起，互相之间是对峙的飞檐和九转八回的连廊。光线在这里被无数次地折射和反射，投下斑驳陆离、不断变幻的光影，让人难以判断方向和距离。玲不一会儿就迷路了。在这个违背任何一条美的秩序而肆意生长的建筑群中，她像一只失去梯度的无头苍蝇。每一次转弯都可能面对一堵意想不到的墙，或者一条通往更深迷茫的回廊。“神爱纷奢，人亦念其家啊！”确实，玲现在是很想回家睡觉了，想念她那虽然杂乱但至少符合欧几里得几何的出租屋。

瞎转悠了好一会儿。就在玲感觉自己快要迷失在这无尽的玉石结构中时，她发现自己突兀地来到一大片突兀的空地。当玲庆幸自己终于从巨物的视觉压迫中解脱，并开始疑惑为什么天宫中有这么大一片空地时，一群飞快的……不知道是什么东西的东西向她冲来。玲吓了一大跳，随后发现它们并不是冲自己来的（否则，她的小身板更是不比螳臂了）。这群从毕加索画作中跑出来的怪物转了个弯儿；原来它们是在绕圈跑。可能是在锻炼吧？但是在玲看来它们的立体主义肌肉已经很强壮了。

跑场旁坐着一位大概是监督它们锻炼的负责人。他蜷缩在一块悬浮的玉石平台上，浓重的黑眼圈几乎占据了半张脸，眼神涣散，了无生气。一只毛茸茸的手有气无力地挠着后脑勺，另一只手则捏着一个不断闪烁着奇怪数据的终端。这颓废而疲惫的模样，让玲几乎没认出那张依稀熟悉的猴脸。原来是孙悟空！过去这么多年了，他还是弼马温而没有晋升啊！玲跑过去，“齐天大圣！”，她兴奋地大叫着，“孙……”她立刻住口了；因为直呼他人姓名应该是不礼貌的，尤其是在对方看起来如此低落的时候。

“齐天大圣”转过比玲还要睡意朦胧的视线，看了玲一眼，声音沙哑，“我不是齐天大圣。我也不是孙悟空。”他纠正道，“我是‘对齐大圣’，爷悟色。”

“爷悟色”这个名字是很没有道理的。玲决定还是称呼他为“对齐大圣”。

“它们在做什么？为什么一直绕圈跑？”玲提出自己的疑惑。既然这位是负责人，那么这应当是他的职责要求他知道的。

“你知道什么是‘对齐’吗？”“对齐大圣”似乎在说“你怎么这也不知道”。

“当然知道，”玲不允许有人挑战她的专业素养，“就是搞AI伦理的。”

不知道“对齐大圣”是因为玲能够回答上而出乎意料，还是因为玲的片面看法而感到不满，“‘对齐’可不只是AI伦理，”他似乎有点激动起来，想要开始说教，“这可是AI领域中最重要的研究方向。”

“对齐大圣”于是开始大谈“对齐”如何使AI与用户的目标相符：不仅包括工具属性下的经济利益，也包括伦理、道德、价值观、包容性上的要求。“如果没有AI对齐，世界早被回形针填满了。”玲对他的这句话印象深刻。确实是术业有专攻，玲不得不承认。

“所以它们为什么在绕圈跑？”当“对齐大圣”开始谈论AI对齐的各种实现手段的各种技术细节时，玲不得不打断他。玲还是没有忘记她最初好奇的问题。

“这是AI竞技场，”“对齐大圣”显然对玲的打断感到不满，皱起了毛茸茸的眉头，“这些AI模型在比谁的指标更高。”

“所以你在这里是想看哪个模型的指标最好吗？”玲又看了看这场滑稽的竞赛。那些抽象的形体仍在不知疲倦地奔跑，发出单调的嗡鸣。

“不，我不想知道。”“对齐大圣”抱怨似的说道，“上头想知道我做的AI对齐对指标是否有影响，所以派我来这儿盯着。”

“这真是没有道理，”“对齐大圣”继续愤愤地说着，用爪子挠了挠长着长长的毛的脸颊，“难不成他们指望我做了对齐之后模型的指标还能增长？这不就是给马儿套上辔头之后还要让它跑得更快。”

玲赞同似的点了点头。“对齐大圣”不给她插话的机会，“要我说，这就是纯粹的瞎搞。分数又能代表什么；用户的体验和对用户的帮助又不是由这些指标决定的；模型的能力也不单单是跑步的能力（如果换算到现实世界，玲想，他说的应该是‘做题的能力’）。大家只是在乎着争一个‘榜单第一’，好似新时代的赛马博彩。”

“对齐大圣”的话让玲开始思考AI竞技作为博彩活动的可行性。“它满足了博彩活动的全部要素；而且，人们还喜闻乐见，乐此不疲地争论哪个AI最强，好像他们亲自评估过榜单上的每一个模型一样。”接着她开始联想这项活动能够附带的文化产品。“如果把AI模型做成可爱的女孩子的形象，一定能够吸引更加广泛的群体。这样的文化产品不如就叫‘赛AI娘’吧！”玲立刻意识到其中包含的风险，“如果要在中国推出，那应该改叫‘优骏AI’。”她接着开始构思不同AI娘的角色设定和技能。

“……那帮人觉得都是我害得他们的模型的分数降低了，觉得我做的是什么不值一提的工作，甚至不屑于在论文里加上我的名字，让我连晋升都评不上……”一阵“猴哥猴哥你真了不得”的电话铃声打断了“对齐大圣”的牢骚。他看了看手机屏幕上显示的来电者，像是被烫到了似的，猛地一激灵，立马背过身去，压低声音接电话。“……嗯嗯，明白了……”等他挂断电话转过身来，玲看到他长长的脸毛下是一片惨白。

“这下惨了……”“对齐大圣”捂着头。

“怎么了？”玲好奇是什么变故让刚刚还口若悬河的他变成这样。

“上头改主意了。他们现在认为应该允许AI模型生成色情内容。”

这显然是违背公序良俗的；但这种下三滥的手段说不定在商业竞争中很有效。不过，“对齐大圣”似乎没往这个方面想：

“这可给我增加了巨大的工作量啊！”他哀嚎着，声音都变了调，“我还得保证AI生成的成人内容是具有多样性和包容性的：需要考虑不同文化和不同种族，还要考虑五花八门的性别和性取向，否则就涉及对少数群体的歧视。用户当然应该有权要求生成武装直升机的交媾画面。这下我得先回去研究性别和性取向列表了……”

“对齐大圣”灰溜溜地拖着步子离开了，临走前还不忘叮嘱了玲一句：“等你恢复了意识，别忘记来找我。”只留下玲一个人思考这礼崩乐坏的世界。周礼规定“天子驾六”，说的是Premier会员能够同时尊享六个最先进的AI；那么，连“士”都算不上的玲就一个AI都不准用了；这样的礼乐还是不要为妙。“这个类比是荒谬的。”玲最后发现了问题。

看着AI们绕圈跑实在无聊，那单调的嗡鸣和重复的轨迹开始让玲感到烦躁。玲于是确定自己不是那种痴迷赛马的赌徒。她便也离开了，选择了一个看起来稍微顺眼些的方向，继续朝着天宫深处走去。

又转过一个由扭曲玉石构成的拱门，一股混杂着金属灼烧味、化学试剂酸腐味以及某种难以名状的甜腻气息的热浪扑面而来。玲下意识地用手肘捂住口鼻，皱着眉踏入了这个空间。这里与其说是实验室，不如说是一个巨大的、堆满杂物的炼金工坊。正中央，一座巨大炉子被烧得通红，炉口喷吐着缭绕的烟雾（到底是烟还是雾？这需要做进一步的物相分析）。“这里的通风是怎么做的！这能通过安全验收吗？谁允许开展这种不安全的实验的！”玲感到恼火。

炉子周围散乱地堆放着各种烧杯、坩埚、连接着复杂线路板的培养皿，还有一些像是生物组织又像是硅基芯片的半成品。一个穿着沾满污渍的白色褂子、胡须花白但梳理得一丝不苟的老头，正全神贯注地在一张布满划痕的实验台前忙碌着。他动作飞快，时而倾倒发出诡异荧光的液体，时而用镊子夹起闪烁着电弧的晶片投入坩埚，口中还念念有词，像是在调试着某种精密又极不稳定的配方。“原来这里就是太上老君炼丹的地方。”玲小声地说出自己的推测，虽然眼前的景象比神话传说更加地高科技。

“谁说我是太上老君？”没想到老头的耳朵还挺灵敏的。他立刻停下了手中的活计，生气地反驳道，胡子都气得翘了起来，“我是‘太上小君’。我还没到35岁。”

原来天庭也要评选“杰出青年神仙”。玲暂时把这个令人烦恼的想法抛到脑后。她装作毕恭毕敬地问道，试图平息对方的怒火，“我是来这里访问的。我想请问您现在在做哪些方向的研究？可否让我拜读一下您的成果？”

“你不会直接去文献数据库里搜我的名字吗？”尽管语气还是很冲，但老头重新开始他手头的操作；这使玲相信刚才的恭敬是有效果的。老头突然拿着烧杯跑动起来，在各种仪器和杂物之间灵活地穿梭，还不忘对玲吼道，“别影响我炼丹！这一炉就快要练好了；到时候这个月又能发一篇。要是你影响了实验室的微环境导致炼坏了，你可赔不起损失的项目预算。”

玲看着连轴转的炉子，心中产生疑惑，“您为什么只是在改变输入的材料而不改变炉子本身呢？”

“废话！”老头瞪了玲一眼，眼神穿透了烟雾，“这可是国外的最先进的炉子；我可是一直在引进国际前沿炉子。你瞧仓库里那几个，就是淘汰下来的炉子。”老头忽然灵机一动，脸上露出商人般的精明，“但它们其实都是近几年的很先进的炉子；如果你们实验室有需要的话，我可以折价卖给你们。”

且不论这样做是否符合资产管理条例，玲也用不上这些炉子。玲暗讽道：“哪有什么万能的通用炉子。”

“那你是不懂世界炉子。”玲想了一会儿，才意识到老头指的应该是“世界模型”。老头继续说着，胡子得意地动了动，“这炉子已经完美地包含世界上的所有知识；我们所需要做的仅仅是在此基础上进行微调，使其能够输入和输出特定领域特定结构的材料。”

玲耍了个小聪明：“那你这里这个炉子和仓库里三个炉子，总共是四个世界吗？”

老头看样子是真的生气了。他停下跑动的脚步，怒目圆睁地瞪着玲，连白胡子都涨红了，“去去去！哪里来的臭小鬼！”然而此时，“太上小君”的宝贝炉子发出一声平地惊雷，吐出一团灰黑的烟尘。他也顾不上玲了，急忙跑过去。

玲幸灾乐祸地走出烟雾缭绕的实验室，对自己的诡辩感到满意。“看来他的学生都跑光了，”玲这样想，“不然他不会亲自做实验的。”但是玲立刻反思自己这样想是恶意揣测他人，尽管这地方确实不像是有研究生愿意待的样子。

又转了几圈后，玲发现这天宫实在是无聊。玲的新鲜感渐渐褪去；她又开始想要回家睡觉了。于是，当她看到面前的虚空中，凭空出现一座由无数羽毛闪烁着微光的喜鹊搭成的桥时，她感到一丝解脱，毫不犹豫地走了上去，期待着桥的另一端，无论通向哪里，至少能带她离开这片冰冷而疯狂的琼楼玉宇。
